\begin{abstract}
Multi-turn mathematical reasoning with small language models (SLMs, $\sim$1B--10B parameters) suffers from substantial performance degradation as conversation history accumulates, with recent work documenting up to 39\% accuracy drops over extended interactions. We address this challenge by proposing \textit{Sparse Logic State Compression} (SLSC), a twin-agent framework that maintains reasoning state as atomic mathematical propositions rather than unstructured text. A reasoning agent generates natural-language chains of thought while a curator agent distills accumulated context into structured propositions, replacing raw history with a compressed, addressable state. We evaluate SLSC on Math-100-multi, a 100-problem benchmark derived from MATH-500 where each problem decomposes into 2--5 dependent sub-questions requiring numerical results from prior turns. Across four Qwen3.5 scales (0.8B--9B), we find that compression methods substantially outperform raw history accumulation at 4B and above: SLSC achieves up to 85\% fully correct rate versus 62\% for uncompressed baselines, while reducing per-turn input tokens by 62\%. However, natural-language summary shows comparable or larger gains, suggesting benefits arise primarily from bounded context size rather than structured representation. A cautionary finding emerges from Lean~4 formalization experiments: augmenting propositions with formal type stubs consistently degrades accuracy, attributed to stateless autoformalization without verifier feedback. Our work establishes structured state compression as an effective, training-free strategy for multi-turn mathematical reasoning in SLMs, while revealing that the core mechanism is preventing attention dilution rather than providing typed mathematical structure.
\end{abstract}
